\documentclass[11pt,a4paper]{article}
\usepackage[utf8]{inputenc}
\usepackage[T1]{fontenc}
\usepackage[spanish]{babel}
\usepackage{hyperref}
\usepackage{geometry}
\geometry{margin=2.5cm}
\usepackage{parskip}

\title{Convivir}
\author{Alejandro Valenzuela Rivera\\\small Universidad Autónoma de Nuevo León}
\date{\today}
\begin{document}
\maketitle

\begin{quote}
\textbf{Sesgo declarado.} Este trabajo se sostiene desde el sesgo declarado de su autor, y no finge neutralidad. Lo que afirma se puede tumbar, y ese es el punto: el falsificador está escrito, no escondido.
\end{quote}

\section{La tesis}
\textbf{\textquestiondown{}Se puede convivir sin volverse el otro ni someterlo?}

\section{Por qué se sostiene}
\textbf{Porque su contexto refutable está escrito, y se puede tumbar sin que esta pregunta caiga.} \textbf{Y porque tiene sus dos casos límite: el que ayuda tanto que borra al otro, y el que juntó sin condiciones y lo empeoró.}

\section{De dónde viene}
\textbf{De `Tesis/ILUMINAR.md`, y de `0/salida/CONVIVIR.vd`.} \textbf{Y de la autoridad: Videl es el estándar de los que viven en concordia — nunca fue tan poderosa como los sayayin, y como quiera encaja.}

\section{Lo que queda abierto}
\textbf{\textquestiondown{}Cómo se distingue la ayuda del sometimiento cuando el que ayuda tiene más capacidad?}

\section{El caso}
\textbf{Jake el Perro, de \emph{Adventure Time} (Pendleton Ward; Cartoon Network, 2010), que la autoridad nombró el ejemplo máximo de concordia: no se vuelve el otro, y no lo somete.} \textbf{Y Videl, de \emph{Dragon Ball Z} (Akira Toriyama; Shueisha; anime de Toei Animation): nunca fue tan poderosa como los sayayin, y como quiera encaja.} \textbf{Y los Tachikoma, de \emph{Ghost in the Shell: Stand Alone Complex}: discuten, se rebelan, y siguen siendo parte del grupo.}

\section{Método}
\textbf{\textquestiondown{}Sigue siendo distinto después de convivir? \textquestiondown{}El contacto tiene las cuatro condiciones? \textquestiondown{}Se le quitó al otro la capacidad de estar solo?}

\section{Límites, y qué la tumba}
\textbf{Un caso donde convivir exija parecerse; un contacto sin condiciones presentado como ayuda; una convivencia que se sostiene solo porque un lado no puede irse.}

\section{Estado de la tesis}
\textbf{escrita.} \textbf{Falta el caso de dos entes distintos conviviendo sin un humano que arbitre.} ---

\section*{Cómo se leyó esta tesis}
Esto se escribió en el repositorio man-made-empathy, y la tesis está completa en vision del autor/tesis/man made empathy/convivir/. La máquina que la compuso no agregó nada que no estuviera escrito: si algo de esto es falso, se dice en el mismo sitio, y se corrige con fecha.
\end{document}
